% Discussion: interpretation, comparison, and implications. Instruction only.

[PROCESSING INSTRUCTION (Discussion). Interpret the Results against the
thesis. Organize as four short threads and keep the prose at the level of
a senior author who connects evidence to meaning.

(1) What the two VVUQ results mean, kept distinct. First, VVUQ of the
codebase itself (did v0.1.0 build correctly) was green: 51 of 51 tests,
clean static analysis, all 15 examples exited 0. That is the verification
floor and it held. Second, and more meaningful, the behavior of the VVUQ
gate was green: it accepted 1 of 6 candidates and blocked 5. That
asymmetry is the whole point, because the pipeline produced a complete,
well formed deliverable in under a millisecond and the gate still refused
five variants. State the honest caveat plainly: the run verified the
machinery thoroughly but exercised the content shallowly; the single
accept used inputs where the reference equalled the observed values and
the human review flag was simply set true, so validation was effectively
self referential because no independent gold standard was supplied.

(2) What ``blocked'' means, stated so a reader is not misled. Make
explicit that ``block'' here carries no failure or crash connotation:
nothing in the actual deliverable production was blocked, since all five
real pipeline deliverables completed every stage and wrote their files.
The blocks occurred only in the Section 03 gate demonstration, where bad
candidates were constructed on purpose to show the gate refusing them. A
blocked file still exists on disk; it is simply gated from
production or clinical use until the named failure is fixed. Present the
five block meanings in Table~\ref{tab:discussion_blocks} and dwell on case
6, where a single failed check out of six dropped the verification
fraction below 1.0, which is the intended strictness.

(3) Honest account of missing inputs and how to improve. The most
important point of the paper: several inputs were absent or synthetic,
above all an independent validation reference. Summarize the gaps (no live
model backends, no reachable web search, no PDF extraction, no real trial
or patient data, and the rich inputs/ corpus that exists in the repository
but was not wired into the executed pipeline). Then give the improvement
priorities in order: supply an independent validation reference; wire the
real inputs/ corpus through ingestion; make the simulations genuinely
stochastic with more than 3 runs and confidence intervals; exercise the
live model, network, PDF, and physics paths; fix the multibyte chunking
bug and add a test; and emit machine readable results alongside the prose.

(4) Practical standing and the regulatory opportunity. Compare the
autonomous cloud run with a conventional high-end server (better:
near-zero setup, one integrated loop, systematic assurance beyond the
shipped tests; different: ephemeral and commit driven; worse: no live
heavy backends, no GPU or display). Then give a candid read on how close
to real life this is: the assurance machinery is real and reusable today,
while the clinical content is still a reference until fed real data and an
independent truth source. Close with a respectful, forward looking note on
the FDA and other bodies: the VVUQ gate is built to support credibility
frameworks such as FDA \S VI.B and ASME V\&V 40, and accelerating rigorous
LLM VVUQ is how the United States can lead on patient safety, efficacy, and
trial speed. Do not claim any endorsement.

SOURCE FILES:

\begin{table}[H]
\caption{Discussion source map (abbreviated paths under repo root).}
\label{tab:discussion_sources}
\centering
\begin{tabular}{L{4.8cm} L{8.6cm}}
\toprule
\textbf{Source (under repo root)} & \textbf{Use for} \\
\midrule
  execution/README.md : server comparison & The better, different, worse comparison. \\
  execution/03-vvuq/ : ``Why VVUQ is substantial'' & The asymmetry argument and the self referential caveat. \\
  execution/04-stage1-automation/ : multibyte finding & The found bug and its disposition. \\
  imagegen/08-fda-cost-efficiency-bridge/ *.py & Figure data: cost bridge and credibility scores. \\
  Clinical-AI-Demos 07-humanoid final-paper & Contrast with the prior workflow (cite, do not copy). \\
\bottomrule
\end{tabular}
\end{table}

SOURCE DATA to render as a publication table (meanings are fixed):

\begin{table}[H]
\caption{What each blocked case means in the real world.}
\label{tab:discussion_blocks}
\centering
\begin{tabular}{L{5.6cm} L{7.8cm}}
\toprule
\textbf{Blocked case} & \textbf{Real-world meaning} \\
\midrule
  No human review recorded & A required sign-off is missing; procedural, not numeric. \\
  Run divergence (CV 0.365, escalated) & The three simulations disagree too much to trust any; a human is paged. \\
  Reference disagreement (rel err 0.60) & The output does not match the trusted reference; likely wrong. \\
  Only 2 runs, need 3 & Insufficient sampling to estimate uncertainty at all. \\
  Artifact exceeds 200K cap & A structural or policy violation; 1 failed check of 6 (fraction 0.83). \\
\bottomrule
\end{tabular}
\end{table}

FIGURE from imagegen Python data (use the .py for data, never the .png;
do not use image-instruct): render
08-fda-cost-efficiency-bridge/08-fda-cost-efficiency-bridge.py as an
illustrative relative effort cost bridge (index 100, 82, 66, 54, 40, 30,
clearly labeled not dollars) with the measured credibility bullets
(Verification 81.9, Validation 85.75 against a target 80) and the FDA \S
VI.B and ASME V\&V 40 references.

HOW TO PROCESS: weave the threads into one argument that ends where the
Introduction began (assurance over generation). Reference, do not restate,
Tables~\ref{tab:discussion_blocks} and the figure. Cite \cite{sel2025vvuq},
\cite{asmevv40}, \cite{fda2023credibility}, \cite{fda2026realtime},
\cite{anthropiclifesciences}, \cite{kawchak2026prediction}, and
\cite{trooskens2026compiled}.]
